\chapter{Methodology}
\label{chap:methodik}

This chapter describes the methods used in this research. It covers dataset construction, model architectures, and experimental design. The research investigates multitask learning (MTL) for SAR-based forest monitoring using two different datasets to address data limitations mentioned in the previous chapter. Section~\ref{sec:treesat_chm_dataset} presents the TreeSatAI-CHM dataset, which extends the existing TreeSatAI benchmark with canopy height labels derived from German LiDAR data. Section~\ref{sec:treesat_models} describes the model architectures developed for this dataset. These include MLP and U-Net baselines and a multitask architecture for joint species classification and height regression. Section~\ref{sec:imaestro_dataset} presents the iMAESTRO dataset. Section~\ref{sec:imaestro_models} describes the model architectures for this dataset, which integrate allometric constraints into the multitask learning framework.

\section{TreeSatAI-CHM Dataset}
\label{sec:treesat_chm_dataset}

The TreeSatAI benchmark~\cite{ahlswede_treesatai_2023} provides Sentinel-1 imagery with species labels for German forests. This section describes the technical process of extending the dataset with canopy height labels derived from official German LiDAR elevation products.

The study area covers Lower Saxony, Germany, where TreeSatAI~\cite{ahlswede_treesatai_2023} provides 50{,}381 image patches at \SI{60}{\meter} spatial extent. The benchmark covers 20 European tree species grouped into 15 tree genera, with species labels from forest administration records. TreeSatAI uses a spatially distributed sampling strategy across Lower Saxony, capturing diverse tree species across different forest types. For this research, only the Sentinel-1 component is used to investigate SAR-only multitask learning.

Each Sentinel-1 patch is provided at \SI{10}{\meter} resolution and contains two polarization channels. These are VV (co-polarized) and VH (cross-polarized). A third feature, the VV/VH ratio, is also provided as a polarimetric index that correlates with vegetation structure. The resulting three-band SAR input captures both backscatter intensity and polarimetric contrast.

The dataset construction proceeds in two stages. First, a canopy height model (CHM) is generated from digital terrain and surface models and aligned to the TreeSatAI patch grid. Second, the resulting height data is combined with Sentinel-1 backscatter to create input tensors. For each patch, four input bands are assembled: VV backscatter (dB), VH backscatter (dB), VV/VH ratio, and CHM (metres). The first three bands come from the TreeSatAI Sentinel-1 archive, while the fourth is the reprojected CHM. All bands are stacked into a single four-band GeoTIFF per patch.

\subsection{Canopy Height Model Generation}
\label{subsec:chm_generation}

\subsubsection{Elevation Data Sources}

\paragraph{}

For canopy height derivation, the elevation data comes from the Lower Saxony State Office for Geoinformation and Land Surveying (LGLN). LGLN products are derived from airborne LiDAR surveys, which provide higher vertical accuracy (typically 0.1 to 0.2~m RMSE) \cite{landesamt_fur_geoinformation_und_landesvermessung_niedersachsen_lgln_digitales_2025,landesamt_fur_geoinformation_und_landesvermessung_niedersachsen_lgln_digitales_2025-1} compared to spaceborne alternatives like TanDEM-X (2 to 10~m RMSE) \cite{wessel_accuracy_2018}. The \SI{1}{\meter} spatial resolution of LGLN products preserves fine-scale canopy structure that would be lost in coarser global products. The data are freely available through the LGLN open data portal. The temporal alignment between LGLN surveys (2014--2019) and TreeSatAI Sentinel-1 acquisitions (2017--2019) minimizes discrepancies from forest growth or disturbance.

The LGLN provides two products relevant to CHM derivation:

\begin{itemize}
    \item \textbf{DTM (Digital Terrain Model, DGM in German)}: A \SI{1}{\meter} resolution model of bare-earth elevation, derived from airborne LiDAR surveys conducted across Lower Saxony. The DTM represents ground surface elevation with vegetation and buildings removed through point cloud classification and filtering. The product is generated from multiple LiDAR acquisition campaigns and undergoes quality control procedures to ensure consistent vertical accuracy across the coverage area. The data are distributed in the ETRS89/UTM zone 32N coordinate reference system (EPSG:25832) with elevation values referenced to the German height reference system DHHN2016.
    \item \textbf{DSM (Digital Surface Model, DOM in German)}: A \SI{1}{\meter} resolution model of the top surface, including trees, buildings, and other above-ground objects. The DSM is derived from the same LiDAR point clouds as the DTM, but uses first-return pulses to capture the uppermost surface rather than ground-classified returns. This product preserves the canopy top structure and is particularly suitable for forestry applications where the highest vegetation points are of interest.
\end{itemize}

Both products are distributed as tiled GeoTIFF files covering Lower Saxony, with metadata indices in GeoJSON format that specify tile boundaries and download URLs. The tiles follow a regular grid pattern with consistent spatial extent, facilitating automated processing across large areas.

\subsubsection{Generation Procedure}

The canopy height model is computed as
\begin{equation}
    \operatorname{CHM} = \operatorname{DSM} - \operatorname{DTM} 
\end{equation}
This standard approach derives vegetation height from paired elevation products. The resulting CHM represents above-ground object height relative to terrain, corresponding to canopy top height for forested areas.

\paragraph{1) Tile Selection and Download}

The LGLN tile index is loaded as a GeoJSON file, and a spatial index is constructed to identify tiles intersecting each TreeSatAI patch. Of 49{,}573 total tiles, 6{,}210 intersect with TreeSatAI locations, achieving \SI{100}{\percent} coverage. Tiles are downloaded via HTTP from the LGLN open data portal.

\paragraph{2) Nodata Interpolation}

Nodata pixels in DTM or DSM propagate to the CHM output, creating gaps in the derived height model. These gaps typically arise from water bodies, steep terrain causing LiDAR occlusion, or areas excluded during quality control. To address this, nodata values are interpolated using iterative neighbourhood averaging.

The interpolation strategy balances gap-filling effectiveness with preservation of spatial structure. Simple nearest-neighbour filling would create blocky artefacts. Sophisticated geostatistical methods like kriging would be computationally expensive for the large dataset and could introduce unrealistic smoothing. The iterative neighbourhood averaging approach provides a middle ground. For each nodata pixel, the algorithm examines its eight neighbours and replaces the value with the mean of valid neighbours. This process repeats for five iterations, progressively filling gaps from the edges inward while maintaining local height gradients. The five-iteration limit was chosen empirically to fill typical gap sizes (5 to 10~m) without excessive propagation of edge values into larger voids.

\paragraph{3) CHM Computation and Quality Control}

After interpolation, the CHM is computed by subtracting DTM from DSM. To remove artefacts, negative CHM values are set to zero, and values exceeding the 95th percentile are clipped to that threshold. This removes implausible values from buildings, power lines, or interpolation errors while preserving the main distribution.

\begin{figure}[htbp]
    \centering
    \includegraphics[width=0.8\textwidth]{dtm_dsm_chm.png}
    \caption[Example DTM, DSM, and CHM tiles for forested areas in Lower Saxony]{Example DTM, DSM, and CHM tiles for forested areas in Lower Saxony. The top row shows the Digital Terrain Model (DTM) representing bare-earth elevation, the middle row shows the Digital Surface Model (DSM) capturing the top surface including tree canopies, and the bottom row shows the derived Canopy Height Model (CHM).}
    \label{fig:dtm_dsm_chm}
\end{figure}

\paragraph{4) Reprojection and Alignment}

The CHM must be reprojected to match TreeSatAI Sentinel-1 geometry, which uses the WGS84 geographic coordinate system with \SI{10}{\meter} pixel spacing. This reprojection involves both coordinate transformation (from ETRS89/UTM to WGS84) and resampling (from \SI{1}{\meter} to \SI{10}{\meter} resolution). Most patches (\SI{87.7}{\percent}) fall within a single elevation tile, while \SI{12.3}{\percent} require mosaicking from multiple tiles.

The resampling strategy preserves canopy height information during the \num{10}$\times$ resolution reduction. Simple averaging would underestimate heights by including sub-canopy pixels. Nearest-neighbour resampling would create blocky artefacts. A two-stage approach inspired by PrediTree \cite{debary_preditree_2025} was adopted instead. First, a local maximum filter (\SI{10}{\meter} window) identifies the dominant canopy height within each target pixel. This preserves the tallest vegetation returns that correspond to canopy tops. Second, bilinear resampling smooths the maximum-filtered result. This produces spatially continuous height fields without blocky transitions. The combination retains local height maxima, which is important for canopy height characterization. 

\paragraph{5) Validation Against Global Canopy Height Products}

The derived CHM is validated against two global products: ETH Global Canopy Height~\cite{lang_high-resolution_2023} (\SI{10}{\meter} resolution from Sentinel-2 and GEDI LiDAR) and Meta Global Canopy Height~\cite{tolan_very_2024} (\SI{1}{\meter} resolution from Maxar imagery and airborne LiDAR). Comparisons use random pixel samples from overlapping areas.

\begin{figure}[htbp]
    \centering
    \includegraphics[width=0.49\textwidth]{chm_lgln_meta.png}
    \hfill
    \includegraphics[width=0.49\textwidth]{chm_lgln_eth.png}
    \caption[Validation of derived CHM against global canopy height products]{Validation of derived CHM against global canopy height products. Left: Comparison with Meta Global Canopy Height at \SI{1}{\meter} resolution. Right: Comparison with ETH Global Canopy Height at \SI{10}{\meter} resolution. Four representative tiles show side-by-side comparisons demonstrating spatial alignment and similar height patterns.}
    \label{fig:chm_validation}
\end{figure}

Visual inspection shows that the LGLN-derived CHM matches well with global products in terms of spatial alignment and height patterns. The average differences of several meters in RMSE could be due to timing differences, sensor variations, or actual uncertainty. Global products report RMSE values between 5 to 10~m when compared to airborne LiDAR \cite{lang_high-resolution_2023,tolan_very_2024}, indicating that the observed differences are within expected ranges. The LGLN-derived CHM is used as the height reference, acknowledging that the labels have some inherent uncertainty.

\subsection{Data Preprocessing}
\label{subsec:treesat_preprocessing}

With CHM patches aligned to the TreeSatAI grid, the height data are combined with Sentinel-1 backscatter to create input tensors. For each patch, four input bands are assembled: VV backscatter (dB), VH backscatter (dB), VV/VH ratio, and CHM (metres). The first three bands come from the TreeSatAI Sentinel-1 archive, while the fourth is the reprojected CHM. All bands are stacked into a single four-band GeoTIFF per patch.

\begin{figure}[htbp]
    \centering
    \includegraphics[width=0.9\textwidth]{histograms_train_val_test60m.png}
    \caption[Distribution of input bands across training, validation, and test splits]{Distribution of input bands across training, validation, and test splits. Each row shows histograms for one split (Train, Val, Test), with columns displaying the four input bands (VV, VH, VV/VH, CHM). The distributions are similar across splits, indicating successful stratified sampling.}
    \label{fig:histograms_splits}
\end{figure}

\subsubsection{1) Label Filtering}

Labels are filtered based on area fraction according to the TreeSatAI GitHub scripts. Only genera occupying more than \SI{7}{\percent} of patch area are retained, corresponding to approximately \SI{250}{\square\meter} coverage. This threshold addresses a fundamental challenge in multi-label classification. The challenge is how to handle species that occupy only a few pixels within a patch.

This indicates that very small area fractions often represent isolated trees or edge effects rather than meaningful forest composition, and their SAR signatures may be overwhelmed by the dominant species. The resulting filtered dataset retains multi-label complexity (average \num{2.3} genera per patch) while removing spurious labels.

\subsubsection{2) Data Splitting Strategy}

The original TreeSatAI 9:1 train:test split is retained for benchmark comparability, allowing direct comparison with published results. Within the training dataset, stratified sampling creates a validation set with 8:1:1 split ratio of train/val/test. For each patch, the primary label (genus with largest area fraction) determines the stratification. Within each genus stratification group, one ninth of samples are randomly assigned to validation, ensuring all genera are represented in both training and validation splits. This stratification is critical for rare genera, which prevent random splitting from placing all samples of uncommon species in the training set.

A limitation of this approach is that there is no spatial separation between splits. Because nearby samples share similar environmental conditions, management history, and SAR acquisition geometry, they tend to be more similar than distant samples. This spatial autocorrelation may cause performance estimates to be higher than what would be achieved on truly independent data. Additionally, TreeSatAI patches are extracted with slight spatial overlap between neighbouring locations, which further increases the similarity between samples that may end up in different splits.

Ideally, train, validation, and test sets would come from separate geographic regions. However, this is difficult for TreeSatAI because the patches are scattered across Lower Saxony rather than forming continuous areas. Creating spatial separation would require buffer zones of 5 to 10~km between splits, which would greatly reduce the usable data. Another problem is that rare genera tend to occur in specific areas (for example, \textit{Larix} is found mainly in mountainous regions). Spatial blocking could put all samples of certain species into just one split, making it impossible to train or evaluate on those species. The stratified random split is a practical choice that keeps class balance and includes rare species in all splits, even though it does not ensure spatial independence.

\begin{figure}[htbp]
    \centering
    \includegraphics[width=0.9\textwidth]{class_distribution_60m.png}
    \caption[Distribution of patches across genera for training, validation, and test splits]{Distribution of patches across genera for training, validation, and test splits. The bar chart shows patch counts for each genus across the three splits. Picea and Pinus dominate the dataset, while several genera have fewer than 1000 patches. The stratified sampling preserves class proportions across splits, ensuring all genera are represented in training, validation, and test sets.}
    \label{fig:class_distribution_60m}
\end{figure}

\subsubsection{3) Class Weighting}

Tree species are not evenly distributed in the dataset. Common genera like \textit{Picea} and \textit{Pinus} possess many patches, while rare genera like \textit{Larix} and \textit{Betula} have fewer than 500 patches. This creates a problem for multi-label classification. Without correction, models focus on common classes and ignore rare ones because the loss from rare classes is very small.

To address this, inverse frequency weighting is adopted. For a dataset with $N$ total samples and $n_c$ samples of class $c$, the weight is
\begin{equation}
    w_c = \frac{N}{C \cdot n_c},
\end{equation}
where $C$ is the number of classes. This gives higher importance to rare classes during training. When the model misclassifies a rare genus, it contributes more to the loss than when it misclassifies a common genus. Dividing by $C$ keeps the total of all weights proportional to the dataset size, which prevents the loss from becoming too large.

Alternative weighting schemes were considered. Effective number weighting~\cite{cui_class-balanced_2019} and focal loss~\cite{lin_focal_2017} are popular options, but inverse frequency weighting was chosen because it is simple and easy to interpret. The weight scales linearly with class frequency, which makes the behaviour predictable. More complex schemes add extra hyperparameters that need tuning.

The final preprocessing converts GeoTIFF stacks and labels into arrays for deep learning. For each split, 4-band patches are stacked into tensors of shape $({N}, 4, 6, 6)$. Labels are converted to multi-hot vectors. The resulting arrays, class weights, and genus mappings are stored for model training.

\begin{figure}[htbp]
    \centering
    \includegraphics[width=0.8\textwidth]{sample_grid_60m.png}
    \caption[Example patches from training, validation, and test sets]{Example patches from training, validation, and test sets. Each row shows randomly selected patches from one split, with columns displaying the four input bands (VV, VH, VV/VH, CHM).}
    \label{fig:sample_grid_60m}
\end{figure}

\begin{figure}[htbp]
    \centering
    \includegraphics[width=0.9\textwidth]{chm_histograms_by_class_train_val_test_60m.png}
    \caption[CHM value distributions by genus across all splits]{CHM value distributions by genus across all splits. Each panel shows the CHM histogram for one genus. Different genera exhibit slightly different height patterns, with coniferous genera like Picea and Pinus showing taller canopies, while some broadleaf genera show lower average heights.}
    \label{fig:chm_histograms_by_class}
\end{figure}

\section{Model Architectures for TreeSatAI-CHM}
\label{sec:treesat_models}

This section describes the model architectures developed for the TreeSatAI-CHM dataset. Two categories of baseline models are used to establish reference performance: multilayer perceptrons (MLPs) following the original TreeSatAI benchmark, and U-Net architectures suitable for extension to multitask learning.

The choice of model architectures was guided by three considerations. First, the small patch size ($6 \times 6$ pixels) constrains the architectural depth. Excessive downsampling would eliminate spatial structure entirely. Second, the classification and regression tasks require architectures that can handle both discrete and continuous outputs. Third, the goal of multitask learning requires architectures with clear separation between shared and task-specific components.

MLPs serve as a simple baseline that treats each patch as a feature vector, ignoring spatial structure. This provides a lower bound on performance and enables comparison with the original TreeSatAI benchmark. In contrast, U-Nets preserve spatial information through skip connections and are naturally suited to multitask extension through shared encoders and task-specific decoders. By comparing MLP and U-Net baselines, the experiments can quantify architecture differences in performance for these tasks.

\subsection{Baseline Models}

\subsubsection{Multilayer Perceptron Baseline}

The MLP baseline follows the TreeSatAI benchmark architecture to ensure comparability with published results. \textbf{The classification MLP} receives the flattened input tensor (3 channels $\times 6 \times 6 = 108$ features) and processes it through three hidden layers of 512 units each before producing class logits for 15 genera.

The three-layer architecture with 512 units per layer was chosen based on the TreeSatAI benchmark. This capacity is sufficient for the classification task without overfitting. Each hidden layer follows the pattern Linear~$\rightarrow$~BatchNorm~$\rightarrow$~ReLU~$\rightarrow$~Dropout. Batch normalization stabilizes training by normalizing activations. Dropout provides regularization by randomly disabling units during training. The output layer produces raw logits passed through a sigmoid function for multi-label classification. A threshold of 0.5, also adapted from TreeSatAI benchmark, is applied to convert probabilities to binary predictions.

\textbf{The regression MLP} shares the same hidden layer structure but outputs a vector of length 36, reshaped to a spatial map of shape $(1, 6, 6)$. This design choice maintains the spatial output format required for height prediction while using the same feature extraction pathway as classification. To ensure non-negative height predictions (negative heights are physically meaningless) and smooth gradients near zero, the output passes through a softplus activation function

\begin{equation}
    \operatorname{softplus}(x) = \log\bigl(1 + \exp(x)\bigr).
\end{equation}

\subsubsection{U-Net Baseline}

The U-Net architecture~\cite{ronneberger_u-net_2015} uses skip connections between encoder and decoder stages to preserve spatial information during downsampling. However, for the small patch sizes ($6 \times 6$ pixels), standard U-Net architectures with multiple downsampling stages would eliminate spatial structure entirely. The encoder therefore uses only two downsampling stages. It reduces the input from $6 \times 6$ to $3 \times 3$, then to a $1 \times 1$ bottleneck. This minimal depth preserves spatial structure while allowing abstract representations.

The U-Net implementation uses three building blocks designed for efficiency and stability. \textbf{DoubleConv blocks} apply two consecutive $3 \times 3$ convolutions with batch normalization and ReLU activation. The double convolution increases the receptive field while maintaining spatial resolution. Batch normalization ensures stable gradient flow. \textbf{Down blocks} use strided convolutions (stride~2) for learnable downsampling, allowing the network to learn optimal downsampling filters. \textbf{Up blocks} use transposed convolutions for upsampling. They concatenate features with corresponding encoder skip connections to recover spatial details lost during downsampling.

\textbf{The classification U-Net} follows the encoder pathway to extract hierarchical features. After the bottleneck, the decoder reconstructs a feature map at original resolution. A $1 \times 1$ convolution produces per-pixel class activations. Global average pooling across spatial dimensions yields patch-level logits for 15 genera. This design allows the network to learn spatially-varying class probabilities before aggregating to patch-level predictions. This can capture within-patch species heterogeneity. Dropout is applied after the first upsampling stage to prevent overfitting on the relatively small dataset.

\textbf{The regression U-Net} requires full spatial output to predict height at each pixel location. The encoder progressively compresses spatial dimensions ($6\times 6 \rightarrow 3\times 3 \rightarrow 1\times 1$) while expanding feature depth ($3 \rightarrow B \rightarrow 2B \rightarrow 4B$ channels, where $B$ is the base channel count). As spatial resolution decreases, feature channels increase to maintain representational capacity. The decoder mirrors the encoder, restoring spatial resolution through skip connections. The concatenation at each decoder stage combines contextual information from deeper layers with high-resolution spatial details from the encoder. Deeper layers capture large-scale patterns while the encoder preserves fine-scale structure. This multi-scale fusion is critical for accurate height prediction. Canopy height exhibits both local variation (tree crowns) and broader patterns (stand-level structure). A final $1 \times 1$ convolution produces a single-channel output of shape $(1, 6, 6)$. This is passed through softplus activation to ensure non-negative height predictions.

The baseline models are designed with multitask learning extension in mind. The U-Net architecture is particularly suitable for this purpose because the encoder naturally provides shared representations that can serve multiple decoder heads. In subsequent experiments, a multitask U-Net variant shares the encoder while maintaining separate architectures for classification and regression decoders.

\subsection{Multitask Learning Architecture}

The multitask architecture follows the hard parameter sharing paradigm(shared encoders), which has been shown to provide effective regularisation by forcing the encoder to learn representations useful for multiple tasks~\cite{vandenhende_multi-task_2021}. A shared encoder extracts common feature representations processed by task-specific decoder heads, and enables investigation of whether species classification and height regression benefit from common SAR feature representations. This design choice was motivated by 2 factors. First, hard parameter sharing provides implicit regularization by constraining the encoder to learn representations useful for multiple objectives. This can improve generalization compared to single-task models. Second, sharing the encoder reduces the total parameter count compared to training separate models, improving computational efficiency. 

Alternative multitask architectures were considered, including soft parameter sharing (separate encoders with cross-task regularization) and task-specific adapters. However, hard parameter sharing was chosen for its simplicity and proven effectiveness in computer vision tasks. The key question is whether the two tasks are sufficiently related to benefit from shared representations, which the experiments aim to answer.

\subsubsection{Architecture Design}

The MultiTaskUNet extends the U-Net regressor baseline with an additional classification head. The shared encoder processes input through two downsampling stages ($6\times 6 \rightarrow 3\times 3 \rightarrow 1\times 1$), which learns features that must support both tasks. The classification head applies global average pooling on bottleneck features, that produces a vector passed through a fully connected layer to generate class logits. This design leverages the most abstract representation (the bottleneck) for classification and this thesis proposes that species features are best captured by high-level semantic features.

\begin{figure}[htbp]
    \centering
    \includegraphics[width=0.95\textwidth]{treesat-model.png}
    \caption[TreeSatAI-CHM multitask learning architecture]{TreeSatAI-CHM multitask learning architecture. The model uses a shared U-Net encoder with two downsampling stages to extract features from $6\times 6$ SAR patches. Two task-specific decoder heads predict genus classification (via global average pooling and fully connected layers) and canopy height regression (via upsampling and skip connections). The shared encoder learns representations beneficial for both tasks while task-specific heads allow specialized feature refinement.}
    \label{fig:treesat_mtl_architecture}
\end{figure}

The regression decoder uses transposed convolutions and skip connections to restore spatial resolution ($1\times 1 \rightarrow 3\times 3 \rightarrow 6\times 6$). Then, a final $1\times 1$ convolution produces single-channel height predictions through softplus activation. This asymmetric design makes sense given the different task requirements, that classification needs global context since species features apply to the whole patch and regression needs spatial detail since height varies pixel to pixel. Branching classification from the bottleneck while regression uses the full decoder lets the architecture handle both needs.

\subsubsection{Loss Function Design}

\subparagraph{Loss Normalisation}

Training a multitask model requires combining multiple loss functions. This poses a fundamental challenge of loss scale mismatch, that the classification and regression tasks produce losses with different scales and units. Binary cross-entropy for classification typically ranges from 0.1 to 1.0, while RMSE for regression can range from 1 to 10 meters depending on prediction accuracy. Naively summing these losses would allow the larger-scale loss to dominate training. This would effectively ignore the smaller-scale task. To address this, individual task losses are normalised by their running means computed across batches:
\begin{equation}
    \tilde{L}_{\text{cls}} = \frac{L_{\text{cls}}}{\bar{L}_{\text{cls}}}, \qquad \tilde{L}_{\text{reg}} = \frac{L_{\text{reg}}}{\bar{L}_{\text{reg}}}.
\end{equation}
This normalisation brings both losses to similar unit scale and enables more balanced gradient contributions. The running mean is computed with exponential moving average (momentum 0.9), which provides stable normalisation while adapting to changing loss magnitudes during training.

\subparagraph{Task Weighting Strategies}

Fixed weight combination uses
\begin{equation}
    L_{\text{total}} = w_{\text{cls}} \cdot \tilde{L}_{\text{cls}} + w_{\text{reg}} \cdot \tilde{L}_{\text{reg}},
\end{equation}
where $w_{\text{cls}}$ and $w_{\text{reg}}$ are hyperparameters that control the relative importance of each task. Although this approach is simple and interpretable, it requires manual tuning to find optimal weights. Grid search over weight combinations is computationally expensive and optimal weights may vary across datasets or training stages.

An alternative approach learns task weights automatically by modelling homoscedastic uncertainty~\cite{kendall_multi-task_2017}. The key insight is that task-dependent uncertainty can serve as a principled basis for weighting. Tasks with higher inherent uncertainty should receive lower weights because their predictions are less reliable. The general uncertainty-weighted loss formulation for multiple tasks is:
\begin{equation}
    L_{\text{total}} = \sum_t \left( \frac{1}{2\sigma_t^2} L_t + \frac{1}{2} \log(\sigma_t^2) \right),
\end{equation}
where $L_t$ is the loss for task $t$ and $\sigma_t$ is a learnable parameter representing task uncertainty. The first term weights each task loss inversely proportional to its uncertainty, while the second term prevents trivially increasing uncertainties to minimize the loss. Each task loss is normalised by batch-wise mean loss magnitude before applying uncertainty weighting, ensuring that the learned uncertainties reflect relative task difficulty rather than absolute loss scales.

This formulation allows the model to automatically balance task contributions during training. If one task is inherently more difficult or noisy, the optimiser increases its uncertainty parameter, which reduces its influence on the shared encoder. For the TreeSatAI-CHM two-task case (classification and regression), this becomes:
\begin{equation}
    L_{\text{total}} = \frac{1}{2\sigma_{\text{cls}}^2} \tilde{L}_{\text{cls}} + \frac{1}{2} \log(\sigma_{\text{cls}}^2) + \frac{1}{2\sigma_{\text{reg}}^2} \tilde{L}_{\text{reg}} + \frac{1}{2} \log(\sigma_{\text{reg}}^2).
\end{equation}
The experiments compare both weighting schemes to determine which produces better multitask performance.

\subsubsection{Gradient Analysis}

A critical question in multitask learning is whether the tasks are compatible. Do they benefit from shared representations, or do they compete for encoder capacity? Gradient cosine similarity analysis provides insight into this question, which measures the alignment between task-specific gradients at each layer.

At each encoder layer, gradients from classification and regression losses are computed separately:
\begin{equation}
    \cos(\theta_{\text{layer}}) = \frac{\nabla_{\theta_{\text{layer}}} L_{\text{cls}} \cdot \nabla_{\theta_{\text{layer}}} L_{\text{reg}}}{\left\lVert\nabla_{\theta_{\text{layer}}} L_{\text{cls}}\right\rVert \cdot \left\lVert\nabla_{\theta_{\text{layer}}} L_{\text{reg}}\right\rVert}.
\end{equation}
Positive cosine similarity ($\cos(\theta) > 0$) indicates that both tasks benefit from similar parameter updates at that layer and suggests compatible learning objectives, while negative similarity ($\cos(\theta) < 0$) indicates gradient conflict and optimizing one task would harm the other. Zero similarity suggests independence. 

This analysis is performed at three encoder locations (enc1, enc2, bottleneck). It monitors whether tasks converge toward shared representations or compete for encoder capacity. If similarity is consistently positive across layers and training epochs, it suggests that species classification and height regression share beneficial SAR features. Instead, persistent negative similarity would indicate that the tasks are fundamentally incompatible and may not benefit from joint training. The gradient analysis provides interpretable diagnostics for understanding multitask learning dynamics.

\section{iMAESTRO Dataset}
\label{sec:imaestro_dataset}


\paragraph{Overview}

The dataset covers three European forest regions. These are Bauges (France), Milicz (Poland), and Sneznik (Slovenia), spanning over \SI{100000}{\hectare} with approximately 42 million individual trees from 51 species. The three sites were chosen to represent diverse forest types. Bauges exhibits mixed temperate mountain forests with high structural heterogeneity. Milicz represents lowland pine-oak forests characteristic of Central Europe. Sneznik contains old-growth beech-fir forests typical of the Dinaric Alps. This diversity ensures that models trained on iMAESTRO data encounter varied forest conditions, improving potential for generalization.

All structural data are organised on a \SI{25}{\meter} grid resolution, where each cell represents \SI{0.0625}{\hectare}. This resolution was chosen to match typical forest inventory plot sizes and to provide sufficient spatial detail for SAR analysis while maintaining computational tractability. The raw data for each site consist of tree-level CSV files containing individual tree attributes (cell identifier, species name, tree count, diameter at breast height, and tree height) and raster files encoding cell identifiers.

\subsection{Processing Raw iMAESTRO Data}

\subsubsection{1) Deriving Reference Labels}

Individual tree biomass is estimated using species-specific allometric equations rather than generic models. This choice reflects the substantial variation in wood properties across species. Coniferous species like \textit{Picea abies} (wood density \SI{0.43}{\gram\per\centi\meter\cubed}) produce less biomass per unit volume than hardwoods like \textit{Quercus robur} (wood density \SI{0.74}{\gram\per\centi\meter\cubed})~\cite{zanne_data_2009}. Using species-specific parameters ensures that the biomass labels accurately reflect ecological reality. This enables the multitask model to learn meaningful relationships between SAR backscatter, species composition, and biomass.

The allometric equations follow the general form
\begin{equation}
    \operatorname{Biomass}~(\si{\kilogram}) = a \times \operatorname{DBH}^b \times \operatorname{wood\_density},
\end{equation}
where $a$ and $b$ are species-specific allometric coefficients, DBH is diameter at breast height (\si{\centi\meter}), and wood density is in \si{\gram\per\centi\meter\cubed}. The power-law relationship between DBH and biomass reflects the geometric scaling of tree dimensions. As trees grow taller and wider, their volume (and thus biomass) increases faster than their diameter. The exponent $b$ typically ranges from 2.0 to 2.8 depending on species and growth form.

Parameters were compiled from published literature for 52 European temperate forest species~\cite{forrester_generalized_2017}, prioritizing studies from similar climatic regions to ensure transferability. Wood density values were obtained from the global wood density database~\cite{zanne_data_2009,chave_towards_2009}. For species without available parameters, default values were applied ($a = 0.0673$, $b = 2.5$, wood density $= 0.55$~\si{\gram\per\centi\meter\cubed}), representing average values for temperate broadleaf species. These defaults affect fewer than \SI{5}{\percent} of trees in the dataset, as most abundant species have well-established allometric relationships.

\subsubsection{2) Aggregating Cell-Level Data}

Data are aggregated to the \SI{25}{\meter} grid cell level to match the spatial resolution provided from the original dataset. This resolution represents a compromise between preserving spatial detail and ensuring sufficient samples per pixel for robust statistics. At finer resolutions (e.g., \SI{10}{\meter}), many cells would contain few or no trees, leading to noisy labels. At coarser resolutions (e.g., \SI{100}{\meter}), spatial heterogeneity would be lost, reducing the model's ability to learn fine-scale patterns.

For each cell, total biomass is computed by summing individual tree biomass weighted by count, then converted to metric tons per hectare (t/ha). This area-normalized metric is standard in forestry and enables comparison across different spatial scales. Height statistics use the 95th percentile (height95) as the primary canopy height metric. The 95th percentile was chosen over maximum height because it is less sensitive to outliers (e.g., exceptionally tall emergent trees or measurement errors) while still capturing the dominant canopy structure. This metric has been widely adopted in LiDAR-based forest monitoring and provides consistency with the TreeSatAI-CHM dataset.

Species composition is determined by identifying the dominant genus (highest tree count) within each cell, reducing 51 species to 31 genus-level categories. This aggregation was necessary for several reasons. First, many species have very few samples in the dataset, making species-level classification infeasible. Second, SAR backscatter is more sensitive to structural differences (e.g., coniferous vs. broadleaf, tree size) than to subtle species-level variations within a genus. The dominant genus approach (rather than multi-label classification) was chosen because \SI{25}{\meter} cells typically contain a single dominant genus (\SI{78}{\percent} of cells have $>$\SI{80}{\percent} of trees from one genus). Multi-label classification would add complexity without substantial benefit for this dataset.

\begin{figure}[htbp]
    \centering
    \includegraphics[width=0.49\textwidth]{height_hist_per_genus_across_sites.png}
    \hfill
    \includegraphics[width=0.49\textwidth]{biomass_hist_per_genus_across_sites.png}
    \caption[Distribution of biomass and height values by genus across all three sites]{Distribution of biomass and height values by genus across all three sites. Left: Biomass distributions. Right: Height distributions. \textit{Fagus} and \textit{Abies} dominate the dataset with distinct characteristics. Coniferous genera generally show higher biomass values relative to height due to higher wood density.}
    \label{fig:imaestro_biomass_height_by_genus}
\end{figure}

\begin{figure}[htbp]
    \centering
    \includegraphics[width=0.9\textwidth]{genus_distribution.png}
    \caption[Overall genus distribution in the iMAESTRO dataset]{Overall genus distribution in the iMAESTRO dataset. The bar chart shows the total number of cells for each genus across all three study sites. \textit{Fagus} and \textit{Abies} are the most abundant genera.}
    \label{fig:imaestro_genus_distribution}
\end{figure}

\subsubsection{3) Generating Raster Data}

Cell-level statistics are converted to spatially explicit GeoTIFF rasters: biomass (t/ha), height95 (m), and dominant genus (numerical code). The raster transformation uses site-specific coordinate reference systems with consistent \SI{25}{\meter} pixel size.

When converting tree-level data to gridded rasters, the resulting maps can show abrupt jumps between neighboring cells. This happens because each cell is computed independently from the trees it contains. For example, one cell might have a few very tall trees giving it high average height, while the adjacent cell has shorter trees. Similarly, the genus map can show scattered single-cell patches where one genus appears isolated among different surrounding genera. While real forests do have natural variation, these sharp pixel-to-pixel changes can make it harder for the model to learn smooth spatial patterns from SAR data.

To address this, spatial smoothing is applied to make the maps more spatially coherent while keeping real forest patterns intact. The smoothing is intentionally gentle to avoid erasing true forest structure like edges between different stand types.

For continuous variables (biomass and height95), an outlier detection and correction approach is used. For each pixel, the method looks at its $5\times 5$ neighborhood and calculates the median absolute deviation (MAD), which measures how much values vary locally. The MAD is more reliable than standard deviation because it is not affected by a few extreme values. If a pixel's value differs from its neighborhood median by more than two standard deviations (estimated as $1.4826 \times \text{MAD}$), it is considered an outlier. Instead of fully replacing it with the neighborhood median, the pixel value is adjusted halfway toward the median. This preserves gradual changes across the landscape (like biomass increasing from young to mature forest) and keeps sharp boundaries (like forest edges) while removing isolated spikes or dips.

The $5\times 5$ window size (\SI{125}{\meter} $\times$ \SI{125}{\meter}) was chosen to match the scale of SAR backscatter patterns, which typically span multiple pixels due to the sensor's resolution and speckle characteristics.

For the categorical genus layer, a majority filter is applied. This common technique in spatial analysis replaces each pixel's genus code with the most common genus in its $5\times 5$ neighborhood. The majority filter smooths boundaries between different forest types and removes isolated single-pixel misclassifications, which are common when assigning discrete categories to continuous forest gradients. This approach is well-suited for genus classification because forests naturally show spatial clustering. Trees of the same genus tend to grow together due to similar environmental preferences and management practices.

Only pixels with valid data are modified in both smoothing operations, preserving the original spatial extent and preventing artificial expansion of the forest mask into non-forested areas.

\subsection{Processing Sentinel-1 GRD Data}

\paragraph{SAR Data Selection}

Sentinel-1 Ground Range Detected (GRD) products~\cite{torres_gmes_2012} are downloaded from the Alaska Satellite Facility. For each site, a single-date acquisition from the 2019 growing season is selected with Interferometric Wide swath mode and dual polarisation (VV and VH). The single-date approach was chosen to match the temporal snapshot nature of the iMAESTRO structural data, avoiding complications from seasonal vegetation changes or forest disturbances between acquisitions. Growing season acquisitions were prioritized because vegetation with leaves provides stronger volume scattering from forest canopies, enhancing the SAR sensitivity to forest structure.

\paragraph{Radiometric and Geometric Correction}

The GRD products are processed using ESA SNAP software, which provides standardized processing workflows for Sentinel-1 data. The processing chain applies precise orbit files (available 2--3 weeks post-acquisition) to refine satellite position and velocity, which is essential for accurate geocoding. Radiometric calibration converts digital numbers to sigma nought backscatter coefficients, which represent the radar cross-section per unit area in the slant range plane. This calibration accounts for antenna gain patterns and range spreading loss, ensuring that backscatter values are comparable across different acquisition geometries.

Terrain correction is applied using Range--Doppler orthorectification with SRTM \SI{1}{arcsecond} DEM (approximately \SI{30}{\meter} resolution). This step corrects for geometric distortions caused by topography, which can be significant in mountainous areas like Bauges and Sneznik. The terrain correction also reprojects the data from the native SAR geometry to site-specific coordinate reference systems (Lambert-93 for Bauges, PUWG 1992 for Milicz, D96/TM for Sneznik), ensuring perfect alignment with iMAESTRO structural layers.

\paragraph{Resampling and Logarithmic Scaling}

After terrain correction, SAR backscatter is resampled to the \SI{25}{\meter} grid using bilinear interpolation. This resampling averages approximately $6.25$ input pixels (\SI{10}{\meter} resolution) into each output pixel, reducing speckle noise while preserving spatial patterns. Bilinear interpolation was chosen over nearest-neighbour (which would preserve speckle) or cubic interpolation (which would be computationally expensive) as a pragmatic compromise.

Calibrated sigma nought values are converted to decibels (dB) as
\begin{equation}
    \sigma_{\text{dB}} = 10 \times \log_{10}(\sigma_{\text{linear}}).
\end{equation}
The logarithmic transformation is standard in SAR processing because backscatter values span several orders of magnitude (e.g., \numrange{0.01}{1.0} in linear units), and the dB scale compresses this range into a more manageable interval (e.g., \numrange{-20}{0}~dB) that is approximately normally distributed. The dB scale also has the advantage that multiplicative effects in the linear domain (such as attenuation) become additive in the log domain, simplifying interpretation.

Both VV and VH polarisation channels are processed and stored as separate GeoTIFF files. Pixels with invalid backscatter values (e.g., due to layover, shadow, or water bodies) are assigned a NoData value of $-9999$, enabling proper masking during subsequent patch extraction.

\subsection{Preparing Training Data}

\subsubsection{1) Multi-Channel Data Stacking}

For each site, five raster layers are stacked into a single multi-channel array: (0) VH backscatter (dB), (1) VV backscatter (dB), (2) biomass (t/ha), (3) dominant genus (numerical code), and (4) height95 (m). All layers are verified to have identical spatial dimensions and coordinate reference systems. A global valid data mask is computed to identify pixels where all five channels contain valid values, defining the usable extent of each site.

\subsubsection{2) Patch Extraction with Spatial Blocking}

A spatial blocking strategy is implemented to prevent spatial autocorrelation from inflating performance estimates. Unlike the TreeSatAI-CHM dataset (where spatial blocking was infeasible due to distributed sampling), the iMAESTRO data consist of contiguous rasters covering entire forest landscapes, making spatial blocking both possible and necessary. Without spatial separation between train and test sets, the model could exploit spatial autocorrelation to achieve artificially high performance: neighbouring patches are similar due to shared environmental conditions, and a model that memorizes training patch locations could perform well on nearby test patches without truly learning generalizable features.

Each site is divided into non-overlapping rectangular blocks (4$\times$4 patches per block), and entire blocks are assigned to train, validation, or test sets. With patch size of $64\times 64$ pixels and stride of 32 pixels (\SI{50}{\percent} overlap), each block covers \SI{3.2}{\kilo\meter} $\times$ \SI{3.2}{\kilo\meter} on the ground. This block size was chosen to be large enough to capture meaningful forest landscape patterns (multiple stands, topographic variation) while small enough to provide sufficient blocks for splitting across the three sets.

Within each block, patches are extracted using a sliding window with \SI{50}{\percent} overlap. The overlap increases the number of training samples while maintaining spatial independence between blocks (patches from different blocks never overlap). Patches are retained if at least \SI{20}{\percent} of pixels contain valid data and the dominant genus channel contains at least one valid pixel. The \SI{20}{\percent} threshold excludes patches dominated by non-forest areas (water, clearings) while retaining patches near forest edges where there should be at least partial coverage.

\subsubsection{3) Train--Validation--Test Splitting}

Blocks are randomly shuffled and sequentially assigned to train (\SI{75}{\percent}), validation (\SI{15}{\percent}), and test (\SI{10}{\percent}) sets. The assignment algorithm iteratively selects the split that is furthest below its target proportion and assigns the next block to that split, ensuring that the final proportions closely match the targets. A fixed random seed ensures reproducibility across different runs of the preprocessing pipeline.

The spatial blocking approach ensures that validation and test sets truly represent held-out geographic areas, which is critical for assessing generalisation to new locations. The validation set enables hyperparameter tuning and early stopping without contaminating the test set, while the test set provides an unbiased estimate of generalization performance.

Final training data are exported as NumPy arrays in \texttt{.npy} format. For each split, three files are created: patches (4D array of shape $N\times 5\times 64\times 64$), \texttt{labels\_dominant\_genus} (1D array), and \texttt{sites} (1D array). Biomass and height95 values are embedded in channels~2 and~4 of the patch arrays.

\paragraph{Trade-offs and Limitations}

The random block assignment ensures geographic diversity in each split, that all three splits contain blocks from different parts of each landscape, preventing site-specific biases. However, this approach does not guarantee balanced genus distributions across splits, as some genera may be concentrated in specific regions. Empirical analysis confirmed that all major genera (\textit{Fagus}, \textit{Abies}, \textit{Pinus}, \textit{Picea}, \textit{Quercus}, \textit{Fraxinus}) are represented in all three splits, though proportions vary slightly.

\begin{figure}[htbp]
    \centering
    \includegraphics[width=0.32\textwidth]{bauges_blocks.png}
    \hfill
    \includegraphics[width=0.32\textwidth]{milicz_blocks.png}
    \hfill
    \includegraphics[width=0.32\textwidth]{sneznik_blocks.png}
    \caption[Spatial distribution of training, validation, and test blocks]{Spatial distribution of training, validation, and test blocks across the three study sites. Left: Bauges (France), centre: Milicz (Poland), right: Sneznik (Slovenia). Purple blocks represent training set (\SI{75}{\percent}), blue blocks validation (\SI{15}{\percent}), and yellow blocks test (\SI{10}{\percent}). The block-based splitting ensures complete spatial separation between splits.}
    \label{fig:imaestro_blocks}
\end{figure}

\begin{figure}[htbp]
    \centering
    \includegraphics[width=1\textwidth]{sample_patches.png}
    \caption[Example training patches from iMAESTRO dataset]{Example training patches from iMAESTRO dataset showing all five channels. Each row represents one split (train, val, test), with columns showing VV backscatter, VH backscatter, dominant genus, canopy height, and biomass. SAR channels display typical speckle patterns, while height and biomass maps show strong spatial correlation.}
    \label{fig:imaestro_sample_patches}
\end{figure}

\begin{figure}[htbp]
    \centering
    \includegraphics[width=0.9\textwidth]{genus_split_distribution_patch_level.png}
    \caption[Distribution of dominant genus labels across splits]{Distribution of dominant genus labels across train, validation, and test splits. \textit{Fagus} is the most common genus with over 300 training patches, followed by \textit{Abies} and \textit{Pinus}. The proportional distribution is relatively consistent across splits.}
    \label{fig:imaestro_genus_split}
\end{figure}

\begin{figure}[htbp]
    \centering
    \includegraphics[width=0.7\textwidth]{height_biomass_hist_split.png}
    \caption[Distribution of biomass and height values across splits]{Distribution of biomass and height values across train, validation, and test splits. Both variables exhibit right-skewed distributions typical of forest structural attributes. The similar distribution shapes across splits indicate balanced data splitting.}
    \label{fig:imaestro_biomass_height_split}
\end{figure}

\begin{figure}[htbp]
    \centering
    \includegraphics[width=1\textwidth]{all_sites_rasters.png}
    \caption[Complete preprocessing outputs for all three study sites]{Complete preprocessing outputs for all three study sites showing the five data layers (VV backscatter, VH backscatter, dominant genus, canopy height, and biomass).}
    \label{fig:imaestro_site_rasters}
\end{figure}

\section{Model Architectures for iMAESTRO}
\label{sec:imaestro_models}

The iMAESTRO models differ from TreeSatAI-CHM models in several key aspects. First, TreeSatAI-CHM provides small patches with local features and limited spatial context, indicating shallow architectures, while iMAESTRO provides large patches with rich spatial structure, enabling deeper networks that can capture multi-scale forest patterns. The larger patch size ($64\times 64$ vs. $6\times 6$ pixels) enables deeper U-Net architectures with four downsampling stages instead of two, allowing the network to learn hierarchical features at multiple scales. Second, the iMAESTRO dataset includes three prediction tasks (genus segmentation, height regression, biomass regression) instead of two, requiring a three-headed multitask architecture. Third, the physics-based allometric relationships between height and biomass enables integration of domain knowledge through regularization.

This section describes the model architectures developed for the iMAESTRO dataset. The models use larger patch sizes ($64\times 64$ pixels) compared to TreeSatAI-CHM and integrate physics-based allometric constraints.

\subsection{Single-Task Baseline Models}

Three independent U-Net architectures~\cite{ronneberger_u-net_2015} are implemented for genus segmentation, height regression, and biomass regression. The U-Net architecture consists of a contracting encoder path and an expanding decoder path connected by skip connections. Unlike the TreeSatAI-CHM models (which use only two downsampling stages due to small patch size), the iMAESTRO models use four downsampling stages, progressively reducing spatial resolution from $64\times 64$ to $32\times 32$ to $16\times 16$ to $8\times 8$ pixels. This deeper architecture is justified by the larger input size and enables the network to learn features at multiple scales, that fine-scale texture could be interpreted as individual tree crowns, mid-scale patterns as forest stands, and large-scale context as landscape structure.

The encoder doubles feature channels at each stage, following the standard U-Net design principle of trading spatial resolution for feature depth. Each encoder stage applies two $3\times 3$ convolutions with batch normalization and ReLU activation, increasing the receptive field while maintaining stable gradient flow. Downsampling uses strided $3\times 3$ convolutions with stride~2, allowing the network to learn optimal downsampling filters. The decoder mirrors the encoder, upsampling feature maps through transposed convolutions and concatenating them with corresponding encoder features via skip connections. Spatial dropout is applied before the final output layer to prevent overfitting.

\textbf{The segmentation model} predicts genus labels for each pixel in the $64\times 64$ patch, producing a dense semantic segmentation map. The decoder outputs 128 channels at full resolution, passed through a $1\times 1$ convolution to produce logits for 22 genus classes. This pixel-level prediction is more challenging than the patch-level classification used in TreeSatAI-CHM, as it requires the model to delineate boundaries between different forest types within the patch.

To handle class imbalance (\textit{Fagus} and \textit{Abies} dominate the dataset while many genera have few samples), 16 rare genus classes are excluded from training by assigning them an ignore index. This allows the model to concentrate on the six dominant genera that account for over \SI{90}{\percent} of the dataset. Pixels labeled with ignored classes do not contribute gradients during backpropagation, preventing the model from learning spurious patterns from insufficient data. The model is trained using cross-entropy loss, which is standard for multi-class segmentation.

\textbf{The height and biomass regression models} predict canopy height (height95) and aboveground biomass density (t/ha) for each pixel. Both use the same U-Net architecture with single-channel output through a $1\times 1$ convolution. Unlike the TreeSatAI-CHM regression model (which uses softplus activation to ensure non-negative outputs), the iMAESTRO regression models use linear outputs because the masked loss function already excludes invalid pixels where a forest and non-forest map is provided in the iMAESTRO dataset, and the training data contain only positive values.

The height and biomass regression models predict canopy height (height95) and aboveground biomass density (t/ha) for each pixel. Both use the same U-Net architecture with single-channel output through a $1\times 1$ convolution. Unlike the TreeSatAI-CHM regression model (which uses softplus activation to ensure non-negative outputs), the iMAESTRO regression models use linear outputs because the masked loss function already excludes invalid pixels, and the training data contain only positive values.

Models are trained using masked RMSE loss, which ignores pixels with invalid target values (NoData or outside the forest mask). This masking is critical because the iMAESTRO patches often contain non-forest areas (water, clearings) where height and biomass are undefined. The masked loss ensures that only valid forest pixels contribute to training, preventing the model from learning to predict zero values for non-forest areas, which would artificially inflate performance metrics.

\subsection{Multi-Task Learning Model}

\subsubsection{Architecture Design}

A multi-task U-Net is developed to jointly predict genus segmentation, canopy height, and biomass from Sentinel-1 input. The architecture employs a shared encoder following the same four-stage downsampling structure as baseline models. The key design question is how to share information among three tasks: one discrete (segmentation) and two continuous (regression). The chosen architecture uses hard parameter sharing with three independent decoder branches, allowing each task to learn specialized features while benefiting from shared low-level representations.

From the bottleneck, three independent decoder branches reconstruct task-specific predictions. Each decoder branch has its own upsampling layers, skip connections, and DoubleConv blocks, providing task-specific feature refinement. The segmentation decoder applies higher dropout (0.4) compared to regression decoders (0.2) to account for the higher complexity of multi-class segmentation. Each branch terminates with a $1\times 1$ convolution producing task-specific outputs: logits of 7 genera for segmentation, and single-channel maps for height and biomass.

This design enables investigation of whether the three tasks benefit from shared representations. If genus feature, canopy height, and biomass are all encoded in SAR backscatter patterns, the shared encoder should learn features useful for all tasks. Instead, if the tasks require fundamentally different features, the multitask model may underperform compared to single-task baselines.

\begin{figure}[htbp]
    \centering
    \includegraphics[width=0.95\textwidth]{imaestro-model.png}
    \caption[iMAESTRO multitask learning architecture]{iMAESTRO multitask learning architecture. The model uses a shared U-Net encoder with four downsampling stages to process $64\times 64$ SAR patches. Three independent decoder branches reconstruct task-specific predictions: genus segmentation (7 classes), canopy height regression, and biomass regression. The deeper architecture compared to TreeSatAI-CHM enables multi-scale feature learning suitable for larger patches. Each decoder applies task-specific dropout rates and produces dense pixel-wise predictions.}
    \label{fig:imaestro_mtl_architecture}
\end{figure}

\subsubsection{Uncertainty Loss}

The model employs the learned uncertainty weighting approach introduced in Section~3.2.3 (Equation~3.6). With three tasks producing losses at different scales (cross-entropy for segmentation, RMSE for height and biomass), the uncertainty-based approach automatically balances task contributions. For the iMAESTRO three-task case (genus segmentation, height regression, and biomass regression), the universal formulation from Equation~3.6 applies directly with $t \in \{\text{seg}, \text{height}, \text{bio}\}$, where $\sigma_{\text{seg}}$, $\sigma_{\text{height}}$, and $\sigma_{\text{bio}}$ are learnable parameters representing task uncertainties for segmentation, height, and biomass respectively.

This formulation allows the model to discover the optimal balance among three tasks during training. If one task is inherently more difficult or noisy (e.g., biomass estimation from SAR is challenging due to saturation effects), the optimizer increases its uncertainty parameter, reducing its influence on the shared encoder.

\subsubsection{Allometric Constraint Integration}

Forest biomass and height are physically related through allometric scaling relationships~\cite{chave_improved_2014}, where larger trees have more biomass than smaller trees due to geometric scaling. An allometric regularisation term enforces consistency between predicted height $H$ and biomass $B$ according to the power law $B = \exp(\alpha) \times H^\beta$. Taking logarithms to the allometric equation yields a linear relationship:
\begin{equation}
    \log(B) = \alpha + \beta \times \log(H).
\end{equation}

The allometric constraint provides weak supervision by encouraging physically plausible height-biomass combinations even when individual predictions may be uncertain. If the model predicts high biomass but low height (or vice versa), the allometric loss penalizes this inconsistency. Additionally, it couples the two regression tasks, allowing the model to leverage correlations between height and biomass patterns in the training data. It may also improve generalization to forest conditions not well-represented in the training data by preventing the model from learning spurious relationships that violate ecological principles.

Therefore, the noval method of integrating the allometric constraint could help forest parameter extraction leverage the multitask framework and enforce the physical relationship between height and biomass. And the allometric loss is designed to penalise deviations from the expected relationship:
\begin{equation}
    L_{\text{allom}} = \frac{1}{K} \sum_k \bigl(\log(B_k) - \alpha - \beta \times \log(H_k)\bigr)^2,
\end{equation}
where the sum runs over $K$ valid forest pixels (pixels with positive height and biomass predictions). The logarithmic formulation ensures that relative errors are penalized equally across the range of values, preventing the loss from being dominated by high-biomass pixels.

Parameters $\alpha = 0.0673$ and $\beta = 2.5$ represent average scaling for mixed European forest types, derived from published allometric studies~\cite{chave_improved_2014,forrester_generalized_2017}. The exponent $\beta = 2.5$ reflects the fact that biomass scales faster than height, where a tree twice as tall typically has more than twice the biomass due to increased diameter and crown size. The intercept $\alpha$ accounts for wood density and form factors.

The allometric loss is weighted by $\lambda_{\text{allom}}$ relative to task losses. This small weight ensures that the allometric constraint provides gentle regularization without dominating the optimization. The weight was chosen through experiments empirically to balance between enforcing physical consistency (requiring higher weight), and allowing the model to fit the training data (requiring lower weight). The final loss combines uncertainty-weighted task losses with the allometric constraint:
\begin{equation}
    L_{\text{final}} = L_{\text{total}} + \lambda_{\text{allom}} \times L_{\text{allom}}.
\end{equation}

This formulation enables the model to learn task-specific features while maintaining consistency with physical relationships, potentially improving generalization to forest conditions not well-represented in the training data.
